% Préambule commun des schémas TikZ — charte Liouaï
% Chaque schéma commence par % Préambule commun des schémas TikZ — charte Liouaï
% Chaque schéma commence par % Préambule commun des schémas TikZ — charte Liouaï
% Chaque schéma commence par \input{_preambule} puis \begin{document}.
\documentclass[tikz,border=6pt]{standalone}
\usepackage[T1]{fontenc}
\usepackage[utf8]{inputenc}
\usepackage[scaled=0.95]{helvet}
\renewcommand{\familydefault}{\sfdefault}
\usepackage{amsmath,amssymb}
\usepackage[european]{circuitikz}
\usetikzlibrary{arrows.meta,positioning,calc,angles,quotes,decorations.pathreplacing,backgrounds}

% Couleurs du logo Liouaï Digital Education
\definecolor{navy}{HTML}{243743}
\definecolor{vert}{HTML}{28B78D}
\definecolor{vertfonce}{HTML}{167A5C}
\definecolor{vertclair}{HTML}{EAF7F2}
\definecolor{gris}{HTML}{A1A9AE}
\definecolor{gristexte}{HTML}{5C6870}
\definecolor{fond}{HTML}{F3F5F6}
\definecolor{jaune}{HTML}{E5C800}
\definecolor{jauneclair}{HTML}{FFFBD6}

\tikzset{
  every picture/.style={line cap=round, line join=round, color=navy, text=navy},
  fleche/.style={-{Stealth[length=2.6mm,width=2mm]}, line width=1.1pt},
  vecteur/.style={fleche, line width=1.4pt},
  titre/.style={font=\bfseries\small, text=navy},
  note/.style={font=\footnotesize, text=gristexte},
  boite/.style={draw=navy, line width=0.9pt, rounded corners=5pt, fill=fond, align=center, inner sep=6pt},
}
 puis \begin{document}.
\documentclass[tikz,border=6pt]{standalone}
\usepackage[T1]{fontenc}
\usepackage[utf8]{inputenc}
\usepackage[scaled=0.95]{helvet}
\renewcommand{\familydefault}{\sfdefault}
\usepackage{amsmath,amssymb}
\usepackage[european]{circuitikz}
\usetikzlibrary{arrows.meta,positioning,calc,angles,quotes,decorations.pathreplacing,backgrounds}

% Couleurs du logo Liouaï Digital Education
\definecolor{navy}{HTML}{243743}
\definecolor{vert}{HTML}{28B78D}
\definecolor{vertfonce}{HTML}{167A5C}
\definecolor{vertclair}{HTML}{EAF7F2}
\definecolor{gris}{HTML}{A1A9AE}
\definecolor{gristexte}{HTML}{5C6870}
\definecolor{fond}{HTML}{F3F5F6}
\definecolor{jaune}{HTML}{E5C800}
\definecolor{jauneclair}{HTML}{FFFBD6}

\tikzset{
  every picture/.style={line cap=round, line join=round, color=navy, text=navy},
  fleche/.style={-{Stealth[length=2.6mm,width=2mm]}, line width=1.1pt},
  vecteur/.style={fleche, line width=1.4pt},
  titre/.style={font=\bfseries\small, text=navy},
  note/.style={font=\footnotesize, text=gristexte},
  boite/.style={draw=navy, line width=0.9pt, rounded corners=5pt, fill=fond, align=center, inner sep=6pt},
}
 puis \begin{document}.
\documentclass[tikz,border=6pt]{standalone}
\usepackage[T1]{fontenc}
\usepackage[utf8]{inputenc}
\usepackage[scaled=0.95]{helvet}
\renewcommand{\familydefault}{\sfdefault}
\usepackage{amsmath,amssymb}
\usepackage[european]{circuitikz}
\usetikzlibrary{arrows.meta,positioning,calc,angles,quotes,decorations.pathreplacing,backgrounds}

% Couleurs du logo Liouaï Digital Education
\definecolor{navy}{HTML}{243743}
\definecolor{vert}{HTML}{28B78D}
\definecolor{vertfonce}{HTML}{167A5C}
\definecolor{vertclair}{HTML}{EAF7F2}
\definecolor{gris}{HTML}{A1A9AE}
\definecolor{gristexte}{HTML}{5C6870}
\definecolor{fond}{HTML}{F3F5F6}
\definecolor{jaune}{HTML}{E5C800}
\definecolor{jauneclair}{HTML}{FFFBD6}

\tikzset{
  every picture/.style={line cap=round, line join=round, color=navy, text=navy},
  fleche/.style={-{Stealth[length=2.6mm,width=2mm]}, line width=1.1pt},
  vecteur/.style={fleche, line width=1.4pt},
  titre/.style={font=\bfseries\small, text=navy},
  note/.style={font=\footnotesize, text=gristexte},
  boite/.style={draw=navy, line width=0.9pt, rounded corners=5pt, fill=fond, align=center, inner sep=6pt},
}

\begin{document}
\begin{tikzpicture}[x=0.7cm, y=0.7cm]
% Q : 3 x 4
\foreach \i in {0,...,3} \foreach \j in {0,...,2} \draw[gris, fill=fond] (\i, -\j) rectangle ++(1,-1);
\foreach \i in {0,...,3} \draw[draw=vertfonce, fill=vert!35] (\i, -1) rectangle ++(1,-1);
\node[titre] at (2, 0.9) {\texttt{Q}};
\node[note] at (2, -3.6) {forme \texttt{(3, 4)}};
\node[note, align=center] at (2, -4.5) {3 recettes $\times$ 4 ingrédients};
\node[font=\Large] at (4.7, -1.5) {@};
% prix : 4
\foreach \j in {0,...,3} \draw[gris, fill=fond] (5.5, 0.5-\j) rectangle ++(1,-1);
\foreach \j in {0,...,3} \draw[draw=vertfonce, fill=vert!35] (5.5, 0.5-\j) rectangle ++(1,-1);
\node[titre] at (6, 1.4) {\texttt{prix}};
\node[note] at (6, -4.1) {forme \texttt{(4,)}};
\node[font=\Large] at (7.4, -1.5) {$=$};
% résultat : 3
\foreach \j in {0,...,2} \draw[gris, fill=fond] (8.2, -\j) rectangle ++(1,-1);
\draw[draw=vertfonce, fill=vert] (8.2, -1) rectangle ++(1,-1);
\node[titre] at (8.7, 0.9) {\texttt{coût}};
\node[note] at (8.7, -3.6) {forme \texttt{(3,)}};
\node[note, align=center] at (8.7, -4.5) {un coût\\ par recette};
% règle des dimensions
\node[font=\footnotesize, text=vertfonce, align=center] at (4.6, -5.9) {règle : le nombre de \textbf{colonnes} de \texttt{Q}\\ doit égaler la \textbf{longueur} de \texttt{prix} (ici 4)};
\node[note, align=center] at (4.6, 2.6) {chaque coût = une ligne de \texttt{Q} combinée avec \texttt{prix}};
\end{tikzpicture}
\end{document}
